\section{Results}
\label{sec:results}

\subsection{Search, screening, retrieval, and extraction flow}
\label{sec:results-flow}

The final synthesis included the 40 accessible studies that underwent standardized full-text extraction.  The review flow was: 199 records identified, 118 duplicates removed, 81 unique records screened, 37 records excluded during screening, 44 reports sought for retrieval, 4 reports not retrieved or inaccessible, and 40 reports extracted.  Unretrieved reports were excluded from the evidence synthesis.

\begin{figure}[htbp]
\centering
\includegraphics[width=0.78\linewidth]{prisma_flow}
\caption{Flow of record identification, screening, retrieval, and inclusion for the review synthesis.}
\label{fig:results-prisma-flow}
\end{figure}

The extracted set comprised 21 core studies and 19 supporting studies published from 2013 to 2025.  Counts below therefore refer to the 40 extracted studies unless identified as multi-label claim-domain or V\&V-linkage counts.

\subsection{Study characteristics and quality}
\label{sec:results-characteristics-quality}

The quality checklist scores ranged from 5.0 to 7.5 on the eight-item scale, with a mean of 6.21.  The score distribution was 5.0 ($n=3$), 5.5 ($n=6$), 6.0 ($n=13$), 6.5 ($n=9$), 7.0 ($n=7$), and 7.5 ($n=2$).  Extraction confidence was coded as high for 29 studies and medium for 11 studies.  RQ relevance was uneven: RQ1 was high for 24 studies and partial for 16; RQ2 was high for 3 and partial for 37; and RQ3 was high for 2, partial for 29, and not relevant for 9.  The mostly partial RQ2 pattern reflects that most papers stated MBSE value claims but did not directly measure EPS reliability, reproducibility, or integrability outcomes.

\begin{table}[htbp]
\centering
\caption{Study, quality, and claim-evidence summary.}
\label{tab:results-study-quality-summary}
\begin{tabular}{lr}
\toprule
Metric & n \\
\midrule
Extracted studies & 40 \\
Core studies & 21 \\
Supporting studies & 19 \\
Publication year range & 2013--2025 \\
Claim records & 270 \\
Unsupported claims & 52 \\
Weak claims & 139 \\
Moderate claims & 77 \\
Strong claims & 2 \\
High-confidence extractions & 29 \\
Medium-confidence extractions & 11 \\
\bottomrule
\end{tabular}
\end{table}

Structured limitation flags further contextualize quality.  The dominant limitations were incomplete lifecycle coverage ($n=40$), weak claim--evidence linkage ($n=40$), single-case or limited-scope evidence ($n=40$), no independent validation ($n=40$), prototype or future-work maturity ($n=40$), no physical-test or flight evidence ($n=39$), no baseline or comparator ($n=38$), no quantitative evaluation ($n=35$), unavailable model assets ($n=34$), no EPS-specific evidence ($n=33$), interoperability limitations ($n=29$), and no pre-physical verification model ($n=20$).  These flags were treated as evidence-strength constraints rather than as exclusion criteria, and open limitation text was retained separately from the aggregation flags.

\subsection{Modeling focus and EPS relevance}
\label{sec:results-modeling-focus-eps}

Modeling focus and EPS relevance were synthesized as a single mutually exclusive classification.  Thirty-four studies contained some EPS or power relevance: 19 reported concrete power-budget or electrical-architecture artifacts, 13 were EPS-adjacent, and 2 treated the EPS itself as the primary model scope.  The remaining 6 studies did not model EPS explicitly: 4 were system-level nanosatellite or CubeSat MBSE studies without power-specific artifacts, and 2 contributed nanosatellite MBSE or V\&V evidence outside the EPS domain.

\begin{figure}[htbp]
\centering
\includegraphics[width=0.78\linewidth]{eps_relevance_counts}
\caption{Mutually exclusive EPS relevance categories in the included studies.}
\label{fig:results-eps-relevance}
\end{figure}

\begin{table}[htbp]
\centering
\caption{Modeling focus and EPS relevance categories.}
\label{tab:results-eps-relevance-counts}
\begin{tabular}{lr}
\toprule
EPS relevance category & n \\
\midrule
Power-budget/electrical-architecture & 19 \\
EPS-adjacent & 13 \\
System-level only & 4 \\
EPS-specific & 2 \\
Not EPS-relevant & 2 \\
\bottomrule
\end{tabular}
\end{table}

The 19 power-budget/electrical-architecture studies should therefore be read as studies with concrete power-related evidence, not as 19 EPS-centered MBSE studies.  Most of this evidence was embedded in system-level models or preliminary design studies.  Examples include RAX/CubeSat energy and data simulations \cite{paper-0001,paper-0003,paper-0008}, SysML parametric mass and power-margin checks \cite{paper-0080}, Capella/Arcadia preliminary EPS allocation for 3ColStar \cite{paper-0078}, CNES concurrent electrical-architecture simulations \cite{paper-0053}, and PLM/Dymola requirements-in-the-loop power simulations \cite{paper-0059}.  The two EPS-specific studies were the open-source Papyrus/SysML power-subsystem model with Java-based parameter and requirement-traceability extraction \cite{paper-0026} and the requirements-in-the-loop PLM/Dymola CubeSat case in which EPS behavior and state-of-charge simulation were central \cite{paper-0059}.  Not-EPS-relevant studies still contributed general nanosatellite MBSE or V\&V evidence, such as the EIRSAT-1 communications-subsystem verification demonstrator \cite{paper-0063}.

This synthesis is intentionally stricter than simply counting any paper that mentioned power or included a generic EPS block.  EPS appeared frequently as a system concern, but direct EPS-first MBSE evidence remained sparse even after the refactored extraction separated EPS-specific, EPS-adjacent, and power-budget evidence.

\subsection{MBSE approaches, languages, and toolchains}
\label{sec:results-approaches-tools}

Approach classification was dominated by extended and enhanced MBSE rather than fully integrated digital-engineering environments.  Seventeen studies were coded as Extended MBSE, 11 as Enhanced MBSE, 6 as Integrated MBSE/digital engineering, and 6 as Pure MBSE.  This distribution situates the evidence base but does not by itself establish value or V\&V effectiveness.

\begin{table}[htbp]
\centering
\caption{MBSE approach-classification counts.}
\label{tab:results-approach-counts}
\begin{tabular}{lr}
\toprule
Approach classification & n \\
\midrule
Extended MBSE & 17 \\
Enhanced MBSE & 11 \\
Integrated MBSE/digital engineering & 6 \\
Pure MBSE & 6 \\
\bottomrule
\end{tabular}
\end{table}

Tool and language classification showed a SysML-centered set of studies: SysML appeared in 28 studies, custom tools in 17, Cameo/MagicDraw/CATIA Magic in 14, MATLAB/Simulink and STK in 10 each, open-source packages in 9, ModelCenter and PLM/ALM tools in 6 each, Python in 5, Capella/Arcadia in 4, requirements tools in 4, and OpenMBEE in 3.  The high count for ``other'' tools in the controlled summary reflects named toolchains that were preserved as emergent codes rather than forced into loose categories.  Results-relevant toolchains included Papyrus/Eclipse for the EPS-specific open-source case \cite{paper-0026}, Gaphor/UML stereotypes for open-source preliminary-design formalization \cite{paper-0029}, docs-as-code/M30ML for SeaLion documentation and traceability \cite{paper-0037}, Dymola/Modelica in 3DExperience \cite{paper-0059}, and IDM-CIC/Simu-CIC/Opalis for electrical-architecture concurrent engineering \cite{paper-0053}.

Toolchain use was usually reported as a means to connect requirements, architecture, analysis, and verification artifacts.  The strongest toolchain evidence came from executable or data-flow demonstrations, such as MagicDraw/Cameo--ModelCenter--STK--MATLAB integrations for RAX-style CubeSat analysis \cite{paper-0001,paper-0003}, Capella/ForPlan or Capella/open-source simulation coupling \cite{paper-0033,paper-0043}, and SysML/PLM/Dymola requirements-in-the-loop simulation \cite{paper-0059}.  Pure reference-model studies reported reusable structures and V\&V packages, but generally provided less executed analysis evidence \cite{paper-0015,paper-0077}.

\subsection{Claimed outcomes and value-evidence strength}
\label{sec:results-claims-evidence}

We extracted 270 claim records.  Evidence confidence for those claims was classified as Unsupported ($n=52$), Weak ($n=139$), Moderate ($n=77$), and Strong ($n=2$).  The two strong claims both came from a single concurrent-model-based CubeSat mission-design study that reported a quantitative comparison with a traditional feasibility-study process \cite{paper-0061}.  Outcome domains are multi-label and are now derived from controlled outcome codes only; the 270 claims produced 727 claim-domain memberships.

\begin{figure}[htbp]
\centering
\includegraphics[width=0.92\linewidth]{claim_evidence_by_outcome}
\caption{Claim evidence by value-outcome domain.  Domains are multi-label, so counts exceed the number of unique claim records.}
\label{fig:results-claim-evidence}
\end{figure}

\begin{table*}[htbp]
\centering
\scriptsize
\caption{Claim evidence by value-outcome domain.  Domains are multi-label.}
\label{tab:results-claim-evidence-by-domain}
\begin{tabular}{p{0.47\linewidth}rrrrrr}
\toprule
Outcome domain & Total & Unsup. & Weak & Mod. & Strong & Studies \\
\midrule
Traceability / consistency / requirements quality & 115 & 12 & 62 & 41 & 0 & 37 \\
Integrability / interoperability / lifecycle continuity & 110 & 13 & 70 & 27 & 0 & 35 \\
Reproducibility / reuse / openness & 86 & 22 & 56 & 8 & 0 & 35 \\
Cost / schedule / effort / efficiency / automation & 85 & 23 & 35 & 25 & 2 & 33 \\
Reliability / mission assurance / risk & 85 & 20 & 29 & 36 & 0 & 33 \\
Decision support / analysis / trade-space / feasibility & 78 & 10 & 28 & 40 & 0 & 33 \\
Generality / scalability / tailoring & 78 & 20 & 51 & 7 & 0 & 34 \\
Communication / collaboration / knowledge & 58 & 14 & 38 & 6 & 0 & 34 \\
Other / uncategorized value claim & 22 & 3 & 9 & 10 & 0 & 9 \\
Governance / ROI / project decisions & 10 & 2 & 6 & 1 & 1 & 6 \\
\bottomrule
\end{tabular}
\end{table*}

The most frequent domains were traceability/consistency/requirements quality (115 memberships across 37 studies), integrability/interoperability/lifecycle continuity (110 across 35), reproducibility/reuse/openness (86 across 35), and cost/schedule/effort/automation and reliability/mission-assurance/risk (85 each across 33 studies).  EPS-specific or power-relevant claims were concentrated in integrability/interoperability (62 memberships), traceability/requirements quality (61), decision support (58), cost/schedule/effort/automation (53), reproducibility/reuse (47), and reliability/mission assurance/risk (46).  The explicit ``other'' category reflects the open-coding policy: observations that did not fit controlled outcome labels were preserved rather than reclassified through keywords.

Moderate evidence typically came from executed simulations, model-derived calculations, traceability artifacts, or case-study outputs.  Examples include the RAX model's energy/data interaction and solar-panel/battery sizing analyses \cite{paper-0001,paper-0003}, Papyrus/Java power-subsystem parameter-impact and requirement-fulfillment tables \cite{paper-0026}, power and SoC simulation against requirements in a PLM/Dymola workflow \cite{paper-0059}, worst-case battery survival simulation with a reported minimum state of charge in an electrical-architecture process \cite{paper-0053}, and positive mass/power margins from SysML parametrics \cite{paper-0080}.  Project-value metrics were rarer and case-specific, such as a one-month versus three-month feasibility-study comparison and lower estimated mission cost in one concurrent-design case \cite{paper-0061}, or the reported four-month preliminary-design cycle using Arcadia and weekly sprints \cite{paper-0078}.

Unsupported and weak claims were common where papers asserted reduced cost, reduced schedule, improved reliability, improved learning, or greater reuse without a baseline, metric, reuse trial, stakeholder evaluation, integration test, or flight/operational evidence.  For example, the EPS-specific Papyrus study provided moderate evidence for traceability and requirement-constraint checking, but its cost/time savings remained a future quantitative-benefit study rather than an observed outcome \cite{paper-0026}.  Reliability/mission-assurance/risk claims were more visible after controlled recoding than in the previous synthesis: 85 memberships were coded, with 36 moderate and none strong.  The two strong claims were not EPS-specific reliability claims; both concerned cost/schedule effects in a single concurrent-design feasibility comparison \cite{paper-0061}.

\subsection{Verification and validation linkage}
\label{sec:results-vv-linkage}

For V\&V, each study could receive multiple linkage categories based on controlled EPS/V\&V artifact, pre-physical verification, and integration/test-readiness codes.  EPS/power-specific pre-physical verification appeared in 25 studies; parametric, budget, or constraint checks in 23; executable simulation or scenario analysis before build in 22; test cases, procedures, or V\&V activity packages in 21; requirements-to-model traceability or verification matrices in 19; and external tool or data-flow coupling for V\&V in 11.  Twelve studies had no explicit EPS-focused pre-physical V\&V, 10 linked the model to physical, SIL/HIL, or integration-test readiness, and 3 had only generic or planned V\&V linkage.

\begin{figure}[htbp]
\centering
\includegraphics[width=0.92\linewidth]{vv_linkage_counts}
\caption{Pre-physical V\&V linkage categories.  Categories are multi-label by study.}
\label{fig:results-vv-linkage}
\end{figure}

\begin{table*}[htbp]
\centering
\scriptsize
\caption{Pre-physical V\&V linkage counts.  Categories are multi-label by study.}
\label{tab:results-vv-linkage-counts}
\begin{tabular}{p{0.68\linewidth}r}
\toprule
V\&V linkage category & n \\
\midrule
EPS/power-specific pre-physical verification & 25 \\
Parametric/budget/constraint checks & 23 \\
Executable simulation or scenario analysis before build & 22 \\
Test cases/procedures/V\&V activity packages & 21 \\
Requirements-to-model traceability or verification matrix & 19 \\
No explicit EPS-focused pre-physical V\&V & 12 \\
External tool or data-flow coupling for V\&V & 11 \\
Physical/SIL/HIL/integration-test readiness link & 10 \\
Generic/planned V\&V linkage only & 3 \\
\bottomrule
\end{tabular}
\end{table*}

Together, these categories show that V\&V was usually represented through EPS/power analysis, parametric checks, simulation, and planned or modeled test artifacts.  Explicit links from the model to test execution or integration evidence were still less common than pre-physical analysis, but the controlled recoding separated actual readiness links from generic V\&V language.

Requirements traceability was usually represented through SysML satisfy/verify relations, verification matrices, or tool-generated traceability views.  In the EPS-specific case, requirement constraints and power-budget equations were extracted from Papyrus and evaluated by Java, producing respected/not-respected flags \cite{paper-0026}.  In academic nanosatellite project modeling, requirements were linked to ``Verified By'' test cases and behavior diagrams such as Bus System Activation and LEOP flow \cite{paper-0023}.  In the 3ColStar preliminary design, EPS continuous-power requirements were traced to physical components such as solar panels and batteries \cite{paper-0078}.  In the NSS-1 SysML case, verify relationships connected mass and power requirements to constraint blocks, and positive margins were used as verification evidence \cite{paper-0080}.

Parametric and executable V\&V evidence was strongest when requirement checks were tied to simulation or budget outputs.  RAX/CubeSat studies executed mission scenarios and trade studies involving energy collection, battery capacity, solar-panel sizing, data storage, and downlink \cite{paper-0001,paper-0003,paper-0008}.  Later studies linked pre-physical V\&V to open-source model-to-analysis pipelines \cite{paper-0027,paper-0029,paper-0043}, PLM/Dymola requirements-in-the-loop simulation \cite{paper-0059}, and electrical-architecture energy-balance simulations before hardware development \cite{paper-0053}.  Test-case or automated-test evidence was more limited and often not EPS-specific, as in the requirements modeling/mapping matrix for CubeSat tests \cite{paper-0034} and the EIRSAT-1 communications verification model coupled to Python and ground-station hardware \cite{paper-0063}.  The physical/SIL/HIL/integration-test readiness category ($n=10$) was broader than completed hardware testing: most linkage still stopped at readiness, planned integration, or partial software/hardware interfaces rather than closed-loop EPS acceptance evidence.

\subsection{Gap mapping}
\label{sec:results-gap-mapping}

Each study was mapped to five gap categories adapted from the broader space-MBSE gap frame: practical implementation pathway; toolchain fragmentation/interoperability; lightweight/resource-constrained tailoring; governance/project deliverables/ROI/value; and NewSpace/agile/distributed/rapid iteration.  Classifications were coded as addressed with evidence, addressed conceptually, mentioned only, or not addressed.

\begin{figure}[htbp]
\centering
\includegraphics[width=0.92\linewidth]{gap_mapping_counts}
\caption{Gap-mapping counts for the extracted studies.}
\label{fig:results-gap-mapping}
\end{figure}

\begin{table*}[htbp]
\centering
\scriptsize
\caption{Gap mapping counts.}
\label{tab:results-gap-mapping-counts}
\begin{tabular}{p{0.46\linewidth}rrrr}
\toprule
Gap & Evidence & Conceptual & Mentioned & Not addressed \\
\midrule
Practical implementation pathway & 29 & 11 & 0 & 0 \\
Toolchain fragmentation/interoperability & 19 & 13 & 6 & 2 \\
Lightweight/resource-constrained tailoring & 15 & 21 & 3 & 1 \\
Governance/project deliverables/ROI/value & 7 & 20 & 13 & 0 \\
NewSpace/agile/distributed/rapid iteration & 6 & 11 & 15 & 8 \\
\bottomrule
\end{tabular}
\end{table*}

Implementation pathways were the most consistently addressed gap: all 40 studies addressed the pathway at least conceptually, with 29 evidence-backed and 11 conceptual classifications.  Evidence-backed examples included concrete modeling/data-flow steps for EPS traceability \cite{paper-0026}, concurrent electrical-architecture simulation workflow \cite{paper-0053}, and Capella/Arcadia plus agile preliminary-design execution \cite{paper-0078}.  Toolchain fragmentation/interoperability was also frequently addressed, with 19 evidence-backed cases, 13 conceptual cases, 6 mentioned-only cases, and 2 not-addressed cases.  Evidence-backed toolchain examples included executable RAX integrations \cite{paper-0001,paper-0003}, Papyrus-to-Java EPS analysis \cite{paper-0026}, open-source model-to-simulation coupling \cite{paper-0043}, and model-to-test/tool integration in EIRSAT-1 \cite{paper-0063}.

Lightweight/resource-constrained tailoring was addressed with evidence in 15 studies, addressed conceptually in 21, mentioned only in 3, and not addressed in 1.  Evidence-backed tailoring appeared in studies with rapid analytical design-space exploration \cite{paper-0012}, academic project-management and design tailoring \cite{paper-0023}, open/concurrent preliminary-design tooling \cite{paper-0027}, docs-as-code for a university CubeSat \cite{paper-0037}, causal MBSE user evaluation \cite{paper-0046}, low-cost tradespace modeling \cite{paper-0054}, and agile distributed preliminary design \cite{paper-0078}.

Governance, project deliverables, ROI, and value linkage remained weaker than implementation-pathway and toolchain evidence, but the controlled recoding identified 7 evidence-backed studies, 20 conceptual studies, and 13 mentioned-only studies.  Evidence-backed examples included the 3ColStar study, which reported Capella/Arcadia governance, weekly sprints, requirement/design traceability, and a four-month preliminary design cycle \cite{paper-0078}, as well as studies linking models to project deliverables, cost/schedule comparisons, or requirements-in-loop decision artifacts \cite{paper-0023,paper-0059,paper-0061}.  NewSpace/agile/distributed/rapid-iteration evidence was also limited, with 6 evidence-backed studies, 11 conceptual, 15 mentioned-only, and 8 not addressed.

